\paragraph*{04.05.02.01 Hierarchie der Projektziele definieren und entwickeln}
\kcilabel{para:04.05.02.01_HierarchieDerProjektzieleDefinierenUndEntwickeln}{04.05.02.01}
\label{para:04.05.02.01_HierarchieDerProjektzieleDefinierenUndEntwickeln}

% Task
\par Neben den bestehenden Herausforderungen durch den Parallelbetrieb von \gls{r3} und \gls{s4} erforderte die Integration der \gls{adn} zusätzliche Anpassungen. Diese betrafen nicht nur Technik, sondern auch Zusammenarbeit, Zeitplanung, Abstimmungsprozesse und Umgebungsgestaltung. Ich passte Anforderungen und Ziele des Projekts an, um die zeitkritische Integration der \gls{adn} trotz des bevorstehenden Lebensendes der bisherigen \gls{erp}-Lösung umzusetzen, ohne die übergeordneten Ziele des \gls{s4}-Programms zu kompromittieren.

% Action
\par Die Anbindung der vier Vorsysteme der \gls{adn} an die \gls{erp}-Systeme des \gls{adac} war neuartig; etablierte Prozesse fehlten. Deshalb entwickelte ich ein neues Konzept für die Projektdurchführung, das die speziellen Anforderungen der \gls{adn}-Integration berücksichtigt. Dazu gehörte die Entscheidung, die für das Hauptprojekt \gls{r2h} geplante \gls{aif}-Schnittstellentechnik für die \gls{adn} zu nutzen und ein separates Projekt \gls{f2h} mit eigenem Budget und Ressourcen zu starten, um die Integration gezielt zu fokussieren und Risiken beherrschbar zu machen.

% Result
\par Die Anpassung von Anforderungen, Zielen und Projektdurchführung ermöglichte die erfolgreiche Berücksichtigung der \gls{adn}-Integration. Das Projektmanagementtool erlaubte kontinuierliche Überwachung der Schnittstellenentwicklung und gezielte Anpassungen. Dieser Zuschnitt findet im Hauptprojekt \gls{r2h} für Welle 1.2 ff. Anwendung.

% Reflect
\par Aufgrund des schwachen Mandats konnte ich die \gls{adn}-Anbindung nicht von Anfang an als festen Bestandteil des Projekts durchsetzen. Dennoch sicherte das angepasste Konzept eine erfolgreiche Integration der \gls{adn} ohne größere Verzögerungen oder Probleme. Ich hätte die \gls{adn}-Anbindung früher als festen Projektbestandteil verankern sollen, um eine nahtlosere Integration und frühere Problemlösungen zu ermöglichen.